\documentclass[11pt,a4paper]{article}

\usepackage{amsmath,amssymb}
\usepackage{booktabs}
\usepackage[colorlinks=true,linkcolor=blue,citecolor=blue,urlcolor=blue]{hyperref}
\usepackage{geometry}
\geometry{margin=2.5cm}
\usepackage{microtype}

\newcommand{\Es}{E^{*}}
\newcommand{\Order}{\mathcal{O}}
\newcommand{\norm}[1]{\left\lVert #1 \right\rVert}

\title{A matrix-free $\mathcal{H}^2$/FMM operator for the elastic
half-space flexibility problem\\[2pt]
\large strategy, design choices, and validation}
\author{Hcontact}
\date{2026-06-27}

\begin{document}
\maketitle

\begin{abstract}
This note documents the matrix-free hierarchical operator
(\texttt{backend="h2"}) used by Hcontact to apply the Boussinesq half-space
flexibility matrix without ever forming it. The operator is a black-box
fast multipole method (bbFMM) built on tensor-product Chebyshev
interpolation. We explain why this design is preferred over a dense matrix
or a classical $\mathcal{H}$-matrix with ACA-compressed blocks, describe the
algorithm and the choices that make it efficient for a regular grid with a
translation-invariant kernel, and report measured memory, accuracy, and a
comparison against the FFT solver Tamaas.
\end{abstract}

\section{Problem and discretisation}

We solve normal contact of a rigid rough indenter against an elastic
half-space. The surface displacement under a pressure field is the
convolution
\begin{equation}
  u_z(\mathbf{x}) = \int G(\mathbf{x}-\mathbf{x}')\,p(\mathbf{x}')\,\mathrm{d}A',
  \qquad
  G(r) = \frac{1}{\pi \Es\, r}.
\end{equation}
On a uniform $N_s\times N_s$ grid of square elements of side $h=L/N_s$,
discretising with piecewise-constant pressure gives the influence
matrix $S_{ij}$ equal to the Love~\cite{Love1929} closed-form integral of $G$
over a source element, evaluated at a target element centre. Because the geometry
is a regular grid, $S_{ij}$ depends only on the index offset
$(|i_x-j_x|,|i_y-j_y|)$ -- the operator is \emph{translation invariant}, so
all $N^2$ entries ($N=N_s^2$) are served from one $N_s\times N_s$ lookup
table built with $\Order(N)$ kernel evaluations.

The contact problem is the bound- and equality-constrained quadratic
program $\min \tfrac12 \mathbf{p}^\top S\,\mathbf{p} + \mathbf{p}^\top
\mathbf{g}_0$ s.t.\ $\mathbf{p}\ge 0$, $\mathrm{mean}(\mathbf{p})=\bar p$,
solved by the Polonsky--Keer projected conjugate gradient~\cite{PolKeer1999}
(see \texttt{doc/theory/pcg.tex}). \textbf{Each PCG iteration needs two
applications of $S$.} The operator is therefore used only as a
matrix--vector product; we never need explicit matrix entries, an LU
factorisation, or a stored matrix. This is the key observation that selects
the algorithm.

\section{Why matrix-free: the cost landscape}

\begin{table}[h]
\centering
\begin{tabular}{lll}
\toprule
approach & memory & matvec \\
\midrule
dense matrix                       & $\Order(N^2)$        & $\Order(N^2)$ \\
classical $\mathcal{H}$-matrix (ACA) & $\Order(N\log N)$  & $\Order(N\log N)$ \\
FFT convolution (zero-padded)      & $\Order(N)$          & $\Order(N\log N)$ \\
$\mathcal{H}^2$/FMM (this note)    & $\Order(N)$          & $\Order(N)$ \\
\bottomrule
\end{tabular}
\caption{Asymptotic cost. For $N_s=1024$ ($N\approx10^6$) the dense matrix
alone is $\approx 8.8$\,TB, so it is impossible; the classical
$\mathcal{H}$-matrix needs several GB.}
\end{table}

The previous Hcontact backend was a classical $\mathcal{H}$-matrix~\cite{Hackbusch2015}:
a list of admissible blocks, each compressed independently as $A_{ts}\approx U_{ts}
V_{ts}^\top$ by adaptive cross approximation (ACA). This works, but for a
regular half-space kernel it stores \emph{redundant bases}: many blocks at
the same tree level with the same relative offset are the \emph{same matrix},
yet each is approximated and stored separately. The memory is
$\Order(N\log N)$ with a large constant (e.g.\ $6.2$\,GB at $N_s=512$).

The $\mathcal{H}^2$/FMM design removes this redundancy in two ways:
(i)~it stores a \emph{single basis per cluster} rather than factors per
block, replacing $A_{ts}\approx U_{ts}V_{ts}^\top$ with
$A_{ts}\approx V_t\, S_{ts}\, V_s^\top$ where $V_t,V_s$ are shared
interpolation operators and $S_{ts}$ is a small coupling matrix; and
(ii)~it exploits translation invariance so that the $V$'s and $S_{ts}$'s are
\emph{cached by $(\text{level},\text{relative offset})$}. The result is
$\Order(N)$ memory with a small constant.

\section{The strategy: black-box FMM with Chebyshev interpolation}

\subsection{Uniform quad-tree}
A balanced quad-tree partitions the grid down to square leaves of side
$\ell$ (\texttt{leaf\_side}, default~$8$, i.e.\ $64$ DOF/leaf). Boxes store
only \emph{index ranges} $[i_{x0},i_{x0}+s)\times[i_{y0},i_{y0}+s)$, never
explicit index lists, so the tree is $\Order(N/\ell^2)$ tiny structs.

\subsection{Near/far split}
A target box interacts \emph{directly} (exact kernel) with leaves in its
\texttt{near\_radius} neighbourhood (default $1$, the $3\times3$ block) and
\emph{approximately} (interpolation) with the standard FMM interaction list
-- children of the parent's neighbours that are not themselves near
neighbours. The interaction list has bounded size at every level, and
together the near field and the per-level interaction lists cover every pair
$(i,j)$ exactly once.

\subsection{Degenerate-kernel (Chebyshev) approximation}
For a well-separated target box $t$ and source box $s$, the kernel is
approximated by interpolating each argument on $q$ Chebyshev nodes
$\{\xi_a\}$ per axis:
\begin{equation}
  G(\mathbf{x},\mathbf{y}) \approx
  \sum_{a}\sum_{b} S_q(\mathbf{x},\boldsymbol{\xi}_a^{t})\,
  G(\boldsymbol{\xi}_a^{t},\boldsymbol{\xi}_b^{s})\,
  S_q(\mathbf{y},\boldsymbol{\xi}_b^{s}),
\end{equation}
with the 1-D interpolation weight
$S_q(x,\xi_m)=\tfrac1q+\tfrac2q\sum_{n=1}^{q-1}T_n(\xi_m)T_n(x)$
($T_n$ Chebyshev polynomials). In 2-D the operators are tensor products
($r=q^2$ nodes per box), applied as two 1-D passes rather than a Kronecker
product. This is exactly the black-box FMM~\cite{FongDarve2009}: only kernel
\emph{evaluations} are needed, never analytic multipole expansions.

\subsection{The six passes}
A matvec $\mathbf{u}=S\,\mathbf{p}$ runs
\[
\text{P2M}\;\to\;\text{M2M}\!\uparrow\;\to\;\text{M2L}\;\to\;
\text{L2L}\!\downarrow\;\to\;\text{L2P}\;\;+\;\;\text{near}.
\]
\textbf{P2M} compresses leaf source values onto Chebyshev coefficients
$M_s = W^\top X_s W$. \textbf{M2M} aggregates children to parents,
$M_{\text{par}}\mathrel{+}=\sum_c R_c M_{\text{child}}$. \textbf{M2L}
translates source to target coefficients along the interaction list,
$L_t\mathrel{+}=\sum_s K_{ts}M_s$. \textbf{L2L} propagates parents to
children, $L_{\text{child}}\mathrel{+}=R_c^\top L_{\text{par}}$. \textbf{L2P}
evaluates $\mathbf{u}_t\mathrel{+}=W L_t W^\top$ at leaf DOFs. The
\textbf{near} pass adds direct contributions of neighbouring leaves. All
passes are OpenMP-parallel over boxes/leaves with disjoint writes.

\subsection{Translation-invariant caching (the decisive choice)}
On a uniform grid with a translation-invariant kernel:
\begin{itemize}
\item the \textbf{M2M/L2L} transfer matrices depend only on the child
quadrant and $q$ -- there are exactly \textbf{four} $q^2\times q^2$ matrices,
reused at every level (boxes are normalised to $[-1,1]^2$ by their geometric
extent, giving exact $1/2$ scaling between levels);
\item the \textbf{M2L} coupling $K_{ts}[a,b]=g(\boldsymbol{\xi}_a^t-
\boldsymbol{\xi}_b^s)$ depends only on $(\text{level},\,\Delta b_x,\,\Delta
b_y)$ and is cached -- a few hundred distinct $q^2\times q^2$ matrices total;
\item the \textbf{near} stencils depend only on the relative leaf offset --
at most $(2\,\text{near\_radius}+1)^2$ distinct $\ell^2\times\ell^2$ dense
blocks.
\end{itemize}
The far-field point kernel is taken to be the \emph{same} Love cell integral
as the near field, $g(\mathbf{d})=\mathrm{love\_uz}(\mathbf{d},h/2,h/2)/(\pi
\Es)$, evaluated at continuous node offsets. Using the Love integral (rather
than the bare $1/r$) makes near and far consistent, so the only far-field
error is the interpolation error.

\section{Complexity and measured memory}

The caches are $\Order(\log N)$ (couplings) and $\Order(1)$ (near stencils);
the per-box coefficient buffers are $\Order(N/\ell^2)$. Total memory is
$\Order(N)$, dominated asymptotically by the buffers at $\approx 13$
bytes/DOF. Table~\ref{tab:mem} (\texttt{bench\_h2\_memory.py}, $q=6$) shows
the bytes/DOF \emph{decreasing} to this bound while a dense matrix explodes.

\begin{table}[h]
\centering
\begin{tabular}{rrrrrr}
\toprule
$N_s$ & $N$ & coupling & buffers & total & dense $N^2$ \\
      &     & [MiB]    & [MiB]   & [MiB] & [MiB] \\
\midrule
64   & 4\,096      & 0.79 & 0.05  & 1.12  & 128 \\
128  & 16\,384     & 1.19 & 0.19  & 1.66  & 2\,048 \\
256  & 65\,536     & 1.58 & 0.75  & 2.61  & 32\,768 \\
512  & 262\,144    & 1.98 & 3.00  & 5.26  & 524\,288 \\
1024 & 1\,048\,576 & 2.37 & 12.0  & 14.65 & 8\,388\,608 \\
2048 & 4\,194\,304 & 2.77 & 48.0  & 51.05 & 134\,217\,728 \\
\bottomrule
\end{tabular}
\caption{H2 operator memory ($q=6$, $\ell=8$, near\_radius~$=1$). The near
stencil cache is a constant $0.28$\,MiB. At $N_s=2048$ the operator is
$51$\,MiB versus $128$\,TiB for the dense matrix.}
\label{tab:mem}
\end{table}

\section{CPU-time scaling to $N_s=16384$}

Table~\ref{tab:cputime} (\texttt{bench\_h2\_cputime.py}, $q=6$, 20-core node)
reports the operator build time, per-matvec wall time, and peak process RSS
for the \texttt{example\_rough\_contact} problem from $N_s=128$ up to
$N_s=16384$ ($N\approx2.7\times10^8$ unknowns), with the full rough-surface
solve included where a converged solution is practical.

\begin{table}[h]
\centering
\begin{tabular}{rrrrrrr}
\toprule
$N_s$ & $N$ & build [s] & matvec [s] & RSS [GiB] & solve [s] & iters \\
\midrule
128   & 16\,384       & 0.007 & 0.0005 & 0.04 & 0.21  & 90 \\
256   & 65\,536       & 0.011 & 0.0013 & 0.05 & 0.80  & 199 \\
512   & 262\,144      & 0.023 & 0.0052 & 0.08 & 3.64  & 241 \\
1024  & 1\,048\,576   & 0.064 & 0.019  & 0.19 & 48.0  & 906 \\
2048  & 4\,194\,304   & 0.245 & 0.098  & 0.26 & ---   & --- \\
4096  & 16\,777\,216  & 0.92  & 0.34   & 0.92 & ---   & --- \\
8192  & 67\,108\,864  & 3.65  & 1.39   & 3.57 & ---   & --- \\
16384 & 268\,435\,456 & 13.98 & 6.06   & 14.2 & ---   & --- \\
\bottomrule
\end{tabular}
\caption{H2 CPU-time study ($q=6$, all cores). Build and per-matvec both scale
$\approx\Order(N)$ (each $4\times$ in $N$ costs $\approx4\times$). The
$2.7\times10^8$-DOF operator builds in $14$\,s, applies in $6$\,s, and fits in
$14$\,GiB; the dense matrix would need $\sim5\times10^{8}$\,GiB. Full solves
are listed to $N_s=1024$: beyond that the projected-CG iteration count grows,
so the total solve time -- not the operator -- becomes the limiting cost, and
only the operator metrics are reported.}
\label{tab:cputime}
\end{table}

Build and matvec scale linearly in $N$, confirming the $\Order(N)$ design; the
peak RSS is $\approx57$ bytes/DOF (the intrinsic caches are $\approx12$
bytes/DOF, the remainder being interaction-list indices and the interpreter).
The operator therefore reaches grids that are entirely out of reach for a
dense or classical $\mathcal{H}$-matrix representation, at second-scale build
and matvec times.

\section{Accuracy}

The interpolation error decreases geometrically with $q$; the near field is
exact. Against the dense Love operator ($N_s=32$, \texttt{tests/test\_h2.cpp}):
relative $L_2$ matvec error is $1.3\times10^{-4}$ at $q=4$ and
$3.2\times10^{-6}$ at $q=6$. The H2 backend, plugged into the same PCG,
reproduces the Hertz contact radius and peak pressure and matches the
$\mathcal{H}$-matrix contact area exactly.

\section{Comparison with the FFT solver (Tamaas)}

For a fully active rectangular grid the operator is a convolution and an FFT
matvec ($\Order(N\log N)$, $\Order(N)$ memory) is the natural reference;
Tamaas implements this via a zero-padded \texttt{dcfft} operator.
Table~\ref{tab:tamaas} (\texttt{compare\_tamaas\_h2.py}) compares the H2
backend ($q=6$) against Tamaas~2.8.1 on the same self-affine surfaces.

\begin{table}[h]
\centering
\begin{tabular}{rrrrr}
\toprule
$N_s$ & frac.\ (Tamaas) & frac.\ (H2) & rel.\ $L_2$ pressure & H2 mem [MiB] \\
\midrule
64  & 0.1265 & 0.1260 & 3.26\% & 1.12 \\
128 & 0.1215 & 0.1201 & 3.17\% & 1.66 \\
256 & 0.1112 & 0.1098 & 2.92\% & 2.61 \\
\bottomrule
\end{tabular}
\caption{H2 vs Tamaas FFT. Contact fractions agree to $\sim10^{-3}$.}
\label{tab:tamaas}
\end{table}

Timing (single thread, \texttt{OMP\_NUM\_THREADS=1}) is reported in
Table~\ref{tab:tamaas-time}. The H2 per-matvec cost is already below the FFT
at every size, and the ratio improves with $N$ (from $0.78$ at $N_s=64$ to
$0.69$ at $N_s=512$) as the $\Order(N)$ apply overtakes the FFT's
$\Order(N\log N)$. End-to-end the H2 solve is $\sim1.5\times$ faster here,
which also reflects a (somewhat) lower iteration count of the projected CG.

\begin{table}[h]
\centering
\begin{tabular}{rrrrrrr}
\toprule
& \multicolumn{3}{c}{per-matvec [ms]} & \multicolumn{3}{c}{end-to-end solve [s]} \\
\cmidrule(lr){2-4}\cmidrule(lr){5-7}
$N_s$ & Tamaas & H2 & H2/T & Tamaas & H2 & H2/T \\
\midrule
64  &  0.36 &  0.28 & 0.78 &  0.029 & 0.017 & 0.58 \\
128 &  1.63 &  1.29 & 0.79 &  0.204 & 0.128 & 0.63 \\
256 &  6.92 &  5.59 & 0.81 &  1.230 & 0.870 & 0.71 \\
512 & 37.30 & 25.89 & 0.69 & 10.319 & 6.822 & 0.66 \\
\bottomrule
\end{tabular}
\caption{Single-thread timing, H2 ($q=6$) vs Tamaas dcfft FFT
(\texttt{compare\_tamaas\_h2.py}). Per-matvec is the same operation
$\mathbf{u}=S\mathbf{p}$ on both sides.}
\label{tab:tamaas-time}
\end{table}

The residual $\sim3\%$ pressure $L_2$ difference is \emph{not} an H2 error:
the H2 operator reproduces the exact Love operator to $3\times10^{-6}$
($q=6$), whereas Tamaas~2.8.1's \texttt{dcfft} influence coefficients
oscillate $\pm2$--$8\%$ from the exact Love values at $1$--$3$ cell
separations (Gibbs-like), and its effective modulus is $2\Es{}^2$ (both
verified against Hertz theory; see the project README). The difference is
therefore dominated by the FFT reference, and it is the same $\approx3.3\%$
seen for the $\mathcal{H}$-matrix backend. Both H2 and FFT are $\Order(N)$ in
memory; H2 is additionally \emph{local}, which leaves room to skip work
outside the contact set (active-set masking), something a global FFT cannot
do without re-padding.

\section{Design choices, summarised}
\begin{itemize}
\item \textbf{bbFMM over analytic FMM}: kernel-agnostic, needs only
evaluations -- robust to the Love near-field form.
\item \textbf{Love far kernel, not $1/r$}: near/far consistency; error is
interpolation-only.
\item \textbf{Geometric box extent}: exact half-scaling so M2M/L2L collapse
to four cached matrices.
\item \textbf{Cache by (level, offset)}: turns $\Order(N\log N)$ stored
factors into $\Order(\log N)$ couplings $+\,\Order(1)$ stencils.
\item \textbf{Defaults} $q=6$, $\ell=8$, near\_radius~$=1$: $\sim10^{-6}$
accuracy at a few MiB; tunable per problem.
\item \textbf{Additive backend}: the dense and $\mathcal{H}$-matrix paths and
all existing tests are untouched; PCG plugs into \texttt{H2Operator::matvec}.
\end{itemize}

\section{Limitations and outlook}
The current operator runs on the full grid (the contact constraint is
handled by the PCG projection). Natural next steps: active-domain
\emph{masking} to skip near/far work outside the contact patch; rectangular
($n_x\neq n_y$) and locally refined grids; single-precision cache storage;
and an FFT backend for the fully-active rectangular case. None of these
change the public API
(\texttt{H2Operator\{build(),matvec(),print\_statistics()\}}).

\bibliographystyle{unsrt}
\bibliography{references}

\end{document}
